\begin{figure}[htbp]
\centering
\begin{tikzpicture}[scale=1.1]
    % Central compartment box
    \draw[thick, fill=UofSCGarnet!25] (0.5,0.5) rectangle (3.5,2.5);
    \node[font=\large] at (2,1.5) {Central};
    \node[font=\small] at (2,0.9) {$V_1, C_1$};

    % Peripheral compartment box
    \draw[thick, fill=UofSCAtlantic!25] (5.5,0.5) rectangle (8.5,2.5);
    \node[font=\large] at (7,1.5) {Peripheral};
    \node[font=\small] at (7,0.9) {$V_2, C_2$};

    % Input arrow
    \draw[->, thick, UofSCHorseshoe] (-1.5,1.5) -- (0.5,1.5) node[left, pos=0, font=\small] {$u(t)$};
    \node[below, font=\tiny] at (-1.5,1.1) {Infusion};

    % Elimination from central compartment (downward)
    \draw[->, thick, UofSCCongaree] (2,0.5) -- (2,-0.7) node[below, font=\small] {$k_{10}C_1$};
    \node[right, font=\tiny] at (2.3,-0.2) {Elimination};

    % Transfer from central to peripheral (forward)
    \draw[->, thick, UofSCAtlantic] (3.5,1.8) -- (5.5,1.8);
    \node[above, midway, font=\small] {$k_{12}$};

    % Transfer from peripheral to central (backward)
    \draw[->, thick, UofSCGarnet] (5.5,1.2) -- (3.5,1.2);
    \node[below, midway, font=\small] {$k_{21}$};

    % Rate constant definitions
    \node[right, font=\tiny] at (9.2,2.2) {$k_{10}$ = elimination rate};
    \node[right, font=\tiny] at (9.2,1.8) {$k_{12}$ = central $\to$ peripheral};
    \node[right, font=\tiny] at (9.2,1.4) {$k_{21}$ = peripheral $\to$ central};
\end{tikzpicture}
\caption{Two-compartment pharmacokinetic model diagram showing drug infusion, inter-compartment transfer, and elimination with rate constants labeled.}
\label{fig:two_compartment}
\end{figure}
